% ══════════════════════════════════════════════════════════════════════════════
% Lecture 16: Ford-Fulkerson Complexity and Applications
% ══════════════════════════════════════════════════════════════════════════════

% ══════════════════════════════════════════════════════════════════════════════
\section*{Complexity of the Ford-Fulkerson Method}
\addcontentsline{toc}{subsection}{Ford-Fulkerson Complexity}
% ══════════════════════════════════════════════════════════════════════════════

% ─────────────────────────────────────────────────────────────────────────────
\subsection{Upper Bound (Integer Capacities)}

\begin{complexity}{Ford-Fulkerson — Naive Bound}{ffnaive}
Assuming integer capacities:
\begin{itemize}[noitemsep]
  \item Finding an augmenting path in $G_f$ takes $O(E)$ (using BFS or DFS).
  \item In each iteration, the flow value increases by at least 1.
  \item \textbf{Total Complexity: $O(E \cdot |f_{\max}|)$} where $|f_{\max}|$ is the 
        maximum flow value.
\end{itemize}
\end{complexity}

This bound is correct but can be very poor if the capacities are large.

% ─────────────────────────────────────────────────────────────────────────────
\subsection{Problematic Case}

Consider the diamond graph with a central edge of capacity 1 and 
capacities $M = 2^{100}$ on the other edges:

\begin{center}
\begin{tikzpicture}[font=\small, >=Stealth,
    node/.style={circle, draw, fill=blue!8, minimum size=0.62cm, inner sep=1pt}]
  \node[node, fill=green!20] (s)  at (0, 1)   {$s$};
  \node[node]                (m)  at (2, 1)   {};
  \node[node, fill=red!20]   (t)  at (4, 1)   {$t$};
  \draw[->, thick] (s) to[bend left=30]  node[above,font=\tiny]{$M$}   (m);
  \draw[->, thick] (s) to[bend right=30] node[below,font=\tiny]{$M$}   (m);
  \draw[->, thick] (m) to[bend left=30]  node[above,font=\tiny]{$M$}   (t);
  \draw[->, thick] (m) to[bend right=30] node[below,font=\tiny]{$M$}   (t);
  \draw[->, thick] (s) -- node[above,font=\tiny,keyred]{$1$}            (m);
\end{tikzpicture}
\end{center}

If Ford-Fulkerson alternately chooses the top path and the bottom path 
by passing through the central edge (capacity 1), the flow increases by 1 
each time. It would take $2M = 2^{101}$ iterations to reach the optimum $|f_{\max}| = 2M$.

\begin{remark}{Irrational Capacities}{irrational}
If capacities are irrational, the Ford-Fulkerson method may 
\textbf{never terminate}: the sequence of flows may converge to a 
sub-optimal value.
\end{remark}

% ─────────────────────────────────────────────────────────────────────────────
\subsection{Variants with Good Complexity}

There are two augmenting path choices that guarantee efficient 
termination for integer capacities:

\begin{center}
\renewcommand{\arraystretch}{1.5}
\begin{tabular}{lcc}
\toprule
\textbf{Augmenting Path Choice} & \textbf{Number of Iterations} & \textbf{Implementation} \\
\midrule
\textbf{Shortest Path} (BFS) & $\le \dfrac{1}{2} \cdot |E| \cdot |V|$ & BFS \\[6pt]
\textbf{Fattest Path} & $\le |E| \cdot \log(|E| \cdot U)$ & Modified Dijkstra \\
\bottomrule
\end{tabular}
\end{center}

where $U$ is the maximum flow value.

\textbf{Shortest Path (Edmonds-Karp algorithm):} we choose an augmenting 
path with the \emph{fewest edges} possible (BFS) at each iteration. 
The total number of iterations is bounded by $O(|V| \cdot |E|)$, which 
gives a global complexity of $O(|V| \cdot |E|^2)$.

\textbf{Fattest Path:} we choose the augmenting path whose 
bottleneck capacity ($c_f(p) = \min_{(u,v)\in p} c_f(u,v)$) is 
\emph{maximal}. This choice eliminates the problematic case described above.

% ══════════════════════════════════════════════════════════════════════════════
\section*{Maximum Flow Applications}
\addcontentsline{toc}{subsection}{Maximum Flow Applications}
% ══════════════════════════════════════════════════════════════════════════════

% ─────────────────────────────────────────────────────────────────────────────
\subsection{Bipartite Matching}

\subsubsection{Problem}

$N$ students apply for $M$ job offers ($M \ge N$). Each student 
receives several offers. Can we find a matching that assigns each student 
to exactly one job for which they were accepted?

\begin{definition}{Bipartite Graph and Matching}{bipartite}
A \textbf{bipartite graph} $G = (L \cup R, E)$ has two disjoint 
vertex sets $L$ (students) and $R$ (jobs); all edges go from 
$L$ to $R$. A \textbf{matching} is a subset of edges such that each 
vertex is incident to at most one edge of the matching.
\end{definition}

\subsubsection{Reduction to Maximum Flow}

\begin{center}
\begin{tikzpicture}[font=\small, >=Stealth,
    node/.style={circle, draw, fill=blue!8, minimum size=0.62cm, inner sep=1pt},
    student/.style={circle, draw, fill=blue!20, minimum size=0.62cm, inner sep=1pt},
    job/.style={circle, draw, fill=green!20, minimum size=0.62cm, inner sep=1pt}]
  \node[node, fill=green!30] (s) at (0, 2)   {$s$};
  \node[student] (l1) at (2, 3.5) {$\ell_1$};
  \node[student] (l2) at (2, 2)   {$\ell_2$};
  \node[student] (l3) at (2, 0.5) {$\ell_3$};
  \node[job]     (r1) at (4, 3.5) {$r_1$};
  \node[job]     (r2) at (4, 2)   {$r_2$};
  \node[job]     (r3) at (4, 0.5) {$r_3$};
  \node[node, fill=red!20]   (t) at (6, 2)   {$t$};

  % s -> students (capacity 1)
  \draw[->, thick, keyblue] (s)  -- node[above,font=\tiny]{1} (l1);
  \draw[->, thick, keyblue] (s)  -- node[above,font=\tiny]{1} (l2);
  \draw[->, thick, keyblue] (s)  -- node[below,font=\tiny]{1} (l3);

  % students -> jobs (capacity 1, preferences)
  \draw[->, thick] (l1) -- node[above,font=\tiny]{1} (r1);
  \draw[->, thick] (l1) -- node[above,font=\tiny]{1} (r2);
  \draw[->, thick] (l2) -- node[above,font=\tiny]{1} (r2);
  \draw[->, thick] (l2) -- node[below,font=\tiny]{1} (r3);
  \draw[->, thick] (l3) -- node[below,font=\tiny]{1} (r3);

  % jobs -> t (capacity 1)
  \draw[->, thick, keyred] (r1) -- node[above,font=\tiny]{1} (t);
  \draw[->, thick, keyred] (r2) -- node[above,font=\tiny]{1} (t);
  \draw[->, thick, keyred] (r3) -- node[below,font=\tiny]{1} (t);

  \node[below=4pt, font=\small] at (3, -0.3)
        {All edges have capacity 1};
\end{tikzpicture}
\end{center}

The construction is:
\begin{itemize}[noitemsep]
  \item Add source $s$ connected to each student (capacity 1).
  \item Add sink $t$ connected from each job (capacity 1).
  \item Direct student$\to$job edges (capacity 1).
  \item Run Ford-Fulkerson.
\end{itemize}

\begin{theorem}{Max Flow = Max Matching}{bipmatch}
The value of the maximum flow in this network is equal to the size of 
the maximum matching in the bipartite graph.
\end{theorem}

\begin{proof}
\textbf{Matching $\Rightarrow$ Flow.} Given a matching of size $k$, 
we construct a flow of value $k$ by setting 1 on the matching edges and 
on the corresponding $s$-student and job-$t$ edges.

\textbf{Flow $\Rightarrow$ Matching.} Since capacities are integers, 
Ford-Fulkerson produces an integer flow (0 or 1 on each edge). 
Flow conservation implies: each student receives at most 1 unit of flow 
from $s$, and each job sends at most 1 unit to $t$. The student$\to$job 
edges with flow 1 thus form a matching of size $|f|$.
\end{proof}

% ─────────────────────────────────────────────────────────────────────────────
\subsection{Edge-Disjoint Paths}

\subsubsection{Problem}

In winter, some roads may be closed. How many distinct routes 
(sharing no edges) exist between Lausanne ($s$) and Geneva airport ($t$)? 
What is the minimum number of roads that must be closed to block 
all traffic from $s$ to $t$?

\begin{definition}{Edge-Disjoint Paths}{edgedisj}
Two paths are \textbf{edge-disjoint} if they share no edge. 
The \textbf{maximum number of edge-disjoint paths} from $s$ to $t$ is the 
largest set of pairwise edge-disjoint $s \leadsto t$ paths.
\end{definition}

\subsubsection{Reduction to Maximum Flow}

We build a flow network from the road graph:
\begin{itemize}[noitemsep]
  \item Assign a capacity of 1 to each road (in each direction).
  \item Remove anti-parallel edges via the dummy vertex gadget $v'$ 
        (seen in Lecture 14).
  \item Run Ford-Fulkerson.
\end{itemize}

\begin{tcolorbox}[colback=lightgreen, colframe=keygreen, fonttitle=\bfseries,
                  title={Flow/Cut Duality for Edge-Disjoint Paths}]
\[
  \text{max-flow} = \text{maximum number of edge-disjoint paths}
\]
\[
  \text{min-cut} = \text{minimum number of roads to close to isolate }s\text{ from }t
\]
\end{tcolorbox}

The minimum cut identifies exactly the minimal set of roads whose 
closure blocks every path from Lausanne to Geneva.

% ─────────────────────────────────────────────────────────────────────────────
\subsection{Lecture Summary}

\begin{tcolorbox}[colback=lightblue, colframe=keyblue, fonttitle=\bfseries,
                  title={Key Points}]
\begin{itemize}[noitemsep]
  \item \textbf{Naive Complexity:} $O(E \cdot |f_{\max}|)$ — can be exponential 
        if capacities are large or irrational.
  \item \textbf{Edmonds-Karp (BFS):} $\le \frac{1}{2} |E| \cdot |V|$ iterations, 
        thus $O(|V| \cdot |E|^2)$ total.
  \item \textbf{Fattest Path:} $\le |E| \cdot \log(|E| \cdot U)$ iterations.
  \item \textbf{Bipartite Matching:} reduction to capacity-1 flow network; 
        max flow $=$ max matching (integrality guaranteed by Ford-Fulkerson).
  \item \textbf{Edge-Disjoint Paths:} capacity 1 on each edge; 
        max flow $=$ number of edge-disjoint paths; 
        min cut $=$ min number of edges to remove.
\end{itemize}
\end{tcolorbox}